% =====================================================
% 题型：立体几何填空题 (q_solid_pyramid_point_distance.tex)
% =====================================================
% =====================================================
% 难度：★★★ (难题)
% =====================================================
\newcommand{\qSolidPyramidPointDistanceStem}{%
  在四棱锥 $P-ABCD$ 中，$PD \perp$ 底面 $ABCD$，$E$ 是 $PC$ 的中点，点 $F$ 在棱 $BP$ 上，且 $EF \perp BP$，四边形 $ABCD$ 为正方形，$PD = CD = 2$．则点 $F$ 到平面 $BDE$ 的距离为 \blank{2cm}．%
}

\newcommand{\qSolidPyramidPointDistanceFigure}[1][1.2]{%
  \begin{tikzpicture}[scale=#1, >=stealth]
    % Coordinates
    \coordinate (A) at (0, 0);
    \coordinate (B) at (2.0, 0);
    \coordinate (C) at (2.7, 0.7);
    \coordinate (D) at (0.7, 0.7);
    \coordinate (P) at (0.7, 2.7);
    
    % Midpoints and points on edges
    \coordinate (E) at ($(P)!0.5!(C)$);
    \coordinate (F) at ($(B)!0.667!(P)$); % BF:FP = 2:1
    
    % Draw visible edges
    \draw[thick] (P) -- (A);
    \draw[thick] (P) -- (B);
    \draw[thick] (P) -- (C);
    \draw[thick] (A) -- (B) -- (C);
    
    % Draw hidden edges
    \draw[thick, dashed] (A) -- (D);
    \draw[thick, dashed] (C) -- (D);
    \draw[thick, dashed] (P) -- (D);
    
    % Draw plane BDE helper lines
    \draw[thick, dashed, blue] (B) -- (D);
    \draw[thick, dashed, blue] (D) -- (E);
    \draw[thick, dashed, blue] (B) -- (E);
    
    % Draw EF line segment
    \draw[thick, dashed, red] (E) -- (F);
    
    \fill (E) circle (1.2pt);
    \fill (F) circle (1.2pt);
    
    \node[below left] at (A) {\small $A$};
    \node[below right] at (B) {\small $B$};
    \node[right] at (C) {\small $C$};
    \node[left] at (D) {\small $D$};
    \node[above] at (P) {\small $P$};
    \node[right] at (E) {\small $E$};
    \node[right] at (F) {\small $F$};
    
    \ifteacher
      % Teacher-only helper lines for geometric distance ratio
      % Intersection of AC and BD
      \coordinate (O) at ($(A)!0.5!(C)$);
      % Midpoint of PB (l meets PB at M)
      \coordinate (M) at ($(P)!0.5!(B)$);
      
      \fill[gray] (O) circle (1.2pt) node[below, gray] {\small $O$};
      \fill[gray] (M) circle (1.2pt) node[above right, gray] {\small $M$};
      
      % Draw diagonal AC (gray dashed)
      \draw[thick, dashed, gray] (A) -- (C);
      % Draw OM line segment (l || PD)
      \draw[thick, dashed, gray] (O) -- (M);
    \fi
  \end{tikzpicture}%
}

\newcommand{\qSolidPyramidPointDistanceAnswer}{$\dfrac{4\sqrt{3}}{9}$}

\newcommand{\qSolidPyramidPointDistanceSolution}{%
  第一步：确定点 $F$ 的位置。
  
  因为 $PD \perp$ 底面 $ABCD$，且 $ABCD$ 为正方形，
  
  所以有 $BC \perp CD$，且 $PD \perp BC$。
  
  因为 $PD \cap CD = D$，所以
  \[ BC \perp \text{平面 } PDC. \]
  
  又 $PC \subset$ 平面 $PDC$，所以 $BC \perp PC$，即 $\triangle PBC$ 为直角三角形。
  
  因为 $PD = CD = 2$，在等腰 Rt$\triangle PDC$ 中，斜边
  \[ PC = \sqrt{PD^2 + CD^2} = 2\sqrt{2}. \]
  
  在 Rt$\triangle PBC$ 中，因为 $BC = 2$，斜边
  \[ BP = \sqrt{BC^2 + PC^2} = \sqrt{2^2 + (2\sqrt{2})^2} = 2\sqrt{3}. \]
  
  因为 $EF \perp BP$，在 Rt$\triangle PBC$ 中，$E$ 是 $PC$ 的中点，
  
  由于 $\triangle PFE \sim \triangle PCB$，可得：
  \[ \frac{PF}{PC} = \frac{PE}{PB} \implies PF = \frac{PE \cdot PC}{PB} = \frac{\sqrt{2} \times 2\sqrt{2}}{2\sqrt{3}} = \frac{2\sqrt{3}}{3}. \]
  
  因为整个线段 $BP = 2\sqrt{3}$，所以
  \[ PF = \frac{1}{3}BP \implies BF = \frac{2}{3}BP. \]
  
  所以，点 $F$ 满足 $BF : BP = 2 : 3$．

  第二步：建立距离比例关系．
  
  因为直线 $BP$ 与平面 $BDE$ 交于点 $B$，且 $F$ 在线段 $BP$ 上，
  
  根据空间几何比例关系，点 $F$ 到平面 $BDE$ 的距离 $d$ 与点 $P$ 到平面 $BDE$ 的距离 $h_P$ 满足：
  \[ d = \frac{BF}{BP} \cdot h_P = \frac{2}{3} h_P. \]
  
  第三步：计算点 $P$ 到平面 $BDE$ 的距离（等体积法）。
  
  考虑三棱锥 $P-BDE$。由等体积法有
  \[ V_{P-BDE} = V_{E-PBD}. \]
  
  因为 $E$ 是 $PC$ 的中点，所以点 $E$ 到平面 $PBD$ 的距离 $h_E$ 等于点 $C$ 到平面 $PBD$ 距离的 $\frac{1}{2}$。
  
  因为 $O \in BD$，且 $CO \perp BD$。
  
  又 $PD \perp$ 底面 $ABCD$，故过点 $O$ 作 $l \parallel PD$，则 $l \perp CO$。
  
  而 $l \subset$ 平面 $PBD$，且 $l \cap BD = O$。
  
  所以
  \[ CO \perp \text{平面 } PBD. \]
  
  故点 $C$ 到平面 $PBD$ 的距离即为半对角线长
  \[ CO = \sqrt{2}. \]
  
  所以，点 $E$ 到平面 $PBD$ 的距离
  \[ h_E = \frac{1}{2}CO = \frac{\sqrt{2}}{2}. \]
  
  在 Rt$\triangle PDB$ 中，因为 $PD = 2$，$BD = 2\sqrt{2}$，
  
  所以其底面积
  \[ S_{\triangle PBD} = \frac{1}{2} BD \cdot PD = 2\sqrt{2}. \]
  
  则三棱锥的体积为：
  \[ V_{P-BDE} = V_{E-PBD} = \frac{1}{3} S_{\triangle PBD} \cdot h_E = \frac{1}{3} \times 2\sqrt{2} \times \frac{\sqrt{2}}{2} = \frac{2}{3}. \]
  
  在 $\triangle BDE$ 中，计算三边长：
  
  因为 $DE$ 是 Rt$\triangle PDC$ 的斜边中线，所以
  \[ DE = \frac{1}{2}PC = \sqrt{2}. \]
  
  在 Rt$\triangle BCE$ 中（其中 $EC = \sqrt{2}$），斜边
  \[ BE = \sqrt{BC^2 + EC^2} = \sqrt{2^2 + 2} = \sqrt{6}. \]
  
  而底面 $BD = 2\sqrt{2}$。
  
  因为
  \[ DE^2 + BE^2 = (\sqrt{2})^2 + (\sqrt{6})^2 = 8 = BD^2, \]
  
  所以 $\triangle BDE$ 是以 $E$ 为直角顶点的直角三角形。
  
  由
  \[ V_{P-BDE} = \frac{1}{3} S_{\triangle BDE} \cdot h_P, \]
  
  得
  \[ \frac{2}{3} = \frac{1}{3} \times \sqrt{3} \times h_P \implies h_P = \frac{2\sqrt{3}}{3}. \]
  
  所以，点 $F$ 到平面 $BDE$ 的距离为：
  \[ d = \frac{2}{3} h_P = \frac{2}{3} \times \frac{2\sqrt{3}}{3} = \frac{4\sqrt{3}}{9}. \]
}

\newcommand{\qSolidPyramidPointDistanceDiagnosis}{%
  本题难点不在单纯计算，而在于两个比例结构的识别：
  其一，$EF\perp BP$ 表示 $F$ 是 $E$ 在 $BP$ 上的垂足，需要在 Rt$\triangle PBC$ 中通过投影关系确定
  $\frac{BF}{BP}=\frac{2}{3}$；
  其二，因 $B\in$ 平面 $BDE$，且 $F,P,B$ 共线，点到平面 $BDE$ 的距离沿直线 $BP$ 成线性比例，
  即 $d(F,\text{平面 }BDE)=\frac{BF}{BP}d(P,\text{平面 }BDE)$。
  常见误区是直接陷入复杂点面距离计算，而没有先利用比例降维。
}

\newcommand{\qSolidPyramidPointDistanceTeachingNotes}{%
  \item \textbf{重难点解析}：
    \begin{itemize}
      \item 本题有两种稳定方法：坐标向量法和空间比例法。坐标法最稳，设
      $D(0,0,0)$，$C(2,0,0)$，$B(2,2,0)$，$P(0,0,2)$，
      可快速求出 $F$ 坐标和平面 $BDE$ 的法向量。
      \item 几何法的关键是先把 $EF\perp BP$ 理解为投影关系。在 Rt$\triangle PBC$ 中，
      $F$ 是 $E$ 到 $BP$ 的垂足，由投影可得 $PF=\frac13BP$，从而
      $BF:BP=2:3$。
      \item 第二个关键是距离比例：因为 $B\in$ 平面 $BDE$，且 $F,P,B$ 共线，
      所以点到平面 $BDE$ 的距离沿 $BP$ 成比例变化，
      $d_F=\frac{BF}{BP}d_P$。
      \item 求 $d_P$ 时使用等体积法：
      $V_{P-BDE}=V_{E-PBD}$。其中 $E$ 是 $PC$ 中点，因此
      $d(E,\text{平面 }PBD)=\frac12 d(C,\text{平面 }PBD)$。
    \end{itemize}
}
